\documentclass[12pt,a4paper]{article}

\usepackage[utf8]{inputenc}
\usepackage[T1]{fontenc}
\usepackage[margin=2.5cm]{geometry}
\usepackage{hyperref}
\usepackage{microtype}
\usepackage{parskip}
\usepackage{lmodern}

\hypersetup{colorlinks=true, urlcolor=blue!60!black}

\pagestyle{empty}

\begin{document}

\begin{flushright}
Kundai Farai Sachikonye \\
Auf der Lichtung 19, 82178 Puchheim \\
\href{https://github.com/fullscreen-triangle}{github.com/fullscreen-triangle} \\
\texttt{kundai.sachikonye@bitspark.com} \\
(+49) 1626386344 \\[6pt]
16 May 2026
\end{flushright}

\vspace{1em}

Leibniz-Institut für Analytische Wissenschaften --- ISAS --- e.V. \\
Bioimaging Research Group \\
Dortmund, Germany \\
Reference: 395\_2026

\vspace{1.5em}

\textbf{Re: PhD Candidate (m/f/d) --- Bioimaging Research Group}

\vspace{1em}

Dear Hiring Committee,

I am applying for the PhD candidate position in the Bioimaging Research Group.
The tools I have built for fluorescence microscopy analysis and mass spectrometry
are publicly deployed and available for anyone to use. Before reading further,
the hiring committee is welcome to try them directly:

\begin{itemize}
  \item Microscopy image analysis (PNG / JPG / TIFF):
    \href{https://hieronymus-omega.vercel.app/observe/microscopy}{hieronymus-omega.vercel.app/observe/microscopy}
  \item Mass spectrometry experiment design and forward simulation:
    \href{https://lavoisier.vercel.app/experiment}{lavoisier.vercel.app/experiment}
\end{itemize}

Both run client-side on GPU (WebGL2) --- no installation, no account.

\textbf{Fluorescence microscopy.}
The Hieronymus platform is a browser-based GPU observation engine for multi-modal
scientific imaging. The microscopy module accepts fluorescence, brightfield, and
phase-contrast images and encodes them into partition coordinates $(n, l, m)$
within a three-dimensional S-entropy space ($S_k + S_t + S_e = 1$), using
adjustable physical parameters including partition depth, inverse temperature, and
inter-modal coupling constant. Behind the browser interface, the full
\href{https://github.com/fullscreen-triangle/helicopter}{Helicopter} framework
resolves structural ambiguity across twelve independent measurement modalities
through sequential constraint exclusion, reducing the configuration space from
${\sim}10^{60}$ possible assignments to a unique structural solution and achieving
effective sub-nanometre resolution beyond what any single optical channel provides
independently. For immune cell infiltration analysis in murine heart and tumour
tissue sections, this means that cell-type assignments that are ambiguous in a
single fluorescence channel become uniquely determined when brightfield, phase
contrast, and spectral data are incorporated.

\textbf{Mass spectrometry.}
The Lavoisier platform at
\href{https://lavoisier.vercel.app/experiment}{lavoisier.vercel.app/experiment}
implements force-free mass spectrometry via partition depth minimisation: a
researcher designs an experiment in the browser, the tool runs a forward simulation
generating a predicted spectral library, and that library is taken into the lab to
guide acquisition. A Rust-based executable is available for local runs on
large-scale datasets. The underlying
\href{https://github.com/fullscreen-triangle/lavoisier}{Lavoisier} framework
combines a numerical MS1/MS2 analysis pipeline with a visual pipeline that
converts spectra into temporal image sequences for CNN analysis, reaching
98.9\,\% feature extraction accuracy against the NIST reference library and
supporting memory-mapped I/O for mzML datasets exceeding 100\,GB.
MALDI MSI is a first-class modality in the framework, directly addressing the
spatial metabolite mapping of heart and tumour tissue that this position involves.

\textbf{Immunological context.}
The immune cell infiltration question at the centre of this position connects
to theoretical work I have developed in parallel.
My manuscript \textit{On the Thermodynamic Consequences of Categorical Completion
on Biological Systems} models MHC binding affinity as a function of the
categorical richness of the parent protein's isoform and post-translational
modification landscape, achieving AUC 0.81 versus NetMHCpan's 0.68 on
2,847 IEDB peptides ($r = 0.73$, $p < 10^{-12}$).
Cardiac proteins (collagen isoforms) and tumour-associated antigens that govern
T-cell and NK-cell infiltration are precisely the proteins the model predicts
as primary immune targets from first principles.
The ISAS group's experimental imaging data represents the direct in\,vivo
measurement of what the model addresses.

\textbf{Wet-lab background.}
At the Universitätsklinikum Hamburg-Eppendorf I built automated pipelines for
proteomics and glycomics characterisation from LC-MS/MS and MALDI-TOF data.
At the University Medical Center Groningen I conducted UPLC-MS/MS shotgun
proteomics for cancer biomarker discovery.
At TUM / Bayerisches Forschungsinstitut I developed mass spectrometry pipelines
covering peak detection, ion annotation, and multi-omics statistical modelling.
These roles give me direct bench experience with the instrument classes
--- QTOF, quadrupole ion trap, MALDI-TOF --- that MALDI MSI tissue studies rely on.

I hold an MSc in Computational Biology (Universität Tübingen) and an MSc in
Pharmaceutical Biotechnology (HAW Hamburg), and I completed doctoral research in
Computational Lipidomics at TUM. I am legally resident in Germany with full work
authorisation and am available immediately. English is native; German is
conversational.

\vspace{1.5em}
Yours sincerely,

\vspace{2.5em}
\textbf{Kundai Farai Sachikonye}

\end{document}
